\section*{Erd\H{o}s Problem \#352: an exact unit-area triangle}
\subsection*{1. Statement and proof scope}
Does there exist an absolute $c>0$ such that every Lebesgue-measurable
$A\subseteq\mathbb R^2$ of measure at least $c$ contains three points
forming a triangle of area exactly $1$?
We reconstruct Erd\H{o}s's known affirmative answer for unbounded
sets of positive measure, using the Lebesgue density theorem explicitly.
We also prove a sufficient bound for connected sets and the disk
obstruction to a sharp constant. The general finite-measure question
remains unresolved.

\subsection*{2. Unbounded sets of positive measure: a complete proof}
Choose a density point of $A$ and translate it to the origin. The
Lebesgue density theorem supplies $r>0$ such that, writing $B_\rho$
for the open disk of radius $\rho$,
\[
 |B_{2r}\setminus A|<\frac{\pi r^2}{16}.            \tag{1}
\]
The density point need not itself belong to $A$. Since $A$ is unbounded,
choose $z\in A$, $D=|z|$, with
\[
 D-r\ge\max(3,2/r).
\]
Let $J$ be rotation through $90$ degrees. For $x\in B_r$, define
\[
 T(x)=x+\frac{2J(x-z)}{|x-z|^2}.
\]
The displacement is at most $2/(D-r)\le r$, so $T(B_r)\subseteq B_{2r}$.
More importantly,
\[
 \det(x-z,T(x)-z)=2,
\]
and hence $z,x,T(x)$ form a triangle of area exactly $1$ for every $x$.
It remains to ensure that both variable vertices belong to $A$.

The map $v\mapsto v/|v|^2$ has derivative of operator norm $|v|^{-2}$.
Consequently every singular value of $DT(x)$ is at least
$1-2/(D-r)^2>1/2$, and
\[
 |\det DT(x)|\ge1/4.                              \tag{2}
\]
The map is also injective on $B_r$. Indeed, if
$v=\rho e^{i\theta}$ in complex coordinates, then
\[
 T(x)-z=v(1+2i/\rho^2),\qquad
 |T(x)-z|=\sqrt{\rho^2+4/\rho^2}.
\]
The last expression is strictly increasing for $\rho>\sqrt2$.
Its value determines $\rho$, after which the argument determines
$\theta$. Here every $\rho=|x-z|$ exceeds $\sqrt2$.
The inverse function theorem and injectivity justify the usual
change-of-variables formula on this image.

By (1)--(2),
\[
 \big|\{x\in B_r:T(x)\notin A\}\big|
 \le4|B_{2r}\setminus A|<\pi r^2/4.
\]
Also $|B_r\setminus A|<\pi r^2/16$. These two exceptional sets do not
cover $B_r$. Choose $x$ outside both. Then $z,x,T(x)\in A$ give the
required triangle. All three vertices are distinct by the nonzero
determinant. This also proves the infinite-measure case, since a
bounded planar set has finite measure.

\subsection*{3. An elementary connected-set bound}
Suppose next that $A$ is bounded, connected, measurable, and has area
$M>4$. Write $D>0$ for its diameter. Choose $p,q\in A$ at distance
$d$ as close to $D$ as desired. Relative to the line $pq$, let $W$ be
the width of the orthogonal projection of $A$ in the normal direction.
The tangential projection has length at most $D$, so containment in
the corresponding rectangle gives $M\le DW$.
Since the line contains points of $A$, at least one side has points
at distances approaching $W/2$ or greater. Thus
\[
 \sup_{x\in A}\operatorname{area}(p,q,x)\ge dW/4\ge dM/(4D).
\]
Taking $d$ close enough to $D$ supplies an actual triangle of area
greater than $1$; attainment of the diameter or width was not assumed.
The set $A^3$ is connected, and its continuous image under the
unsigned area function is an interval. This interval contains $0$
(use a repeated vertex) and a value greater than $1$, so it contains
$1$. This proves the assertion for connected sets of area greater
than $4$, including the already-treated unbounded case.

\subsection*{4. The sharp-constant obstruction from a disk}
In a disk of radius $R$, the largest triangle has area
$3\sqrt3 R^2/4$. To verify the maximization, fix two vertices:
unsigned area is the absolute value of an affine function of the
third, so its maximum over the disk is attained on the circle.
Successively maximizing each vertex reduces the problem to an
inscribed triangle. If its angles are $A,B,C$, its area is
$2R^2\sin A\sin B\sin C$. Concavity of $\log\sin$ on $(0,\pi)$ and
$A+B+C=\pi$ give the equilateral maximum.

Disks of radius less than $2\,3^{-3/4}$ therefore contain no unit-area
triangle, while their areas approach
\[
 \frac{4\pi}{\sqrt{27}}.
\]
Any universal threshold $c$ in the question must be at least this value.
The connected-set proof above gives a weaker sufficient bound and does
not establish the conjectured sharp one.

\subsection*{5. Assessment and references}
The missing case is a bounded measurable set with unrestricted
disconnectedness. A triangle of area greater than $1$ alone does not
yield one of area exactly $1$ without an additional continuity
argument. The proof above does not silently interpolate outside $A$.

Source: \texttt{https://www.erdosproblems.com/352}, checked 2026-09-25,
which attributes the unbounded positive-measure result to Erd\H{o}s
and identifies the density-theorem method.
The earlier GPT Pro 5.2 and Claude Opus 4.5 files were reviewed.
External standard inputs used here are the Lebesgue density theorem,
change of variables for injective smooth maps, and preservation of
connectedness by products and continuous maps; the geometric deductions
are proved above. No novelty or formal verification is claimed.
\subsection*{COMPLETION ESTIMATE}
\textbf{COMPLETION: 50\%}
